\section{部署模式}

\app\ 支持两种传输路径:\role{LAN 直连}和\role{Relay 中继}。两者在主持人手机端完全无感 —— 同一个网页,扫不同的码而已。

\subsection{LAN 模式 —— 默认}

\textbf{典型场景}:演出现场后台机器和主持人手机在同一个 WiFi(漫展场地路由器,或临时部署一台便携路由器)。

\begin{center}
\begin{tikzpicture}[
  node distance=1.4cm,
  box/.style={rectangle, draw=sdink, line width=0.7pt, rounded corners=2pt, minimum width=2.6cm, minimum height=0.95cm, align=center, font=\small\sffamily},
  flow/.style={-{Stealth[length=4pt]}, line width=0.7pt}
]
  \node[box] (host) {导播工作机\\\scriptsize\color{sddim}Windows · 8787 端口};
  \node[box, right=of host, fill=sdpanel] (rt) {WiFi 路由器};
  \node[box, right=of rt] (phone) {主持人手机\\\scriptsize\color{sddim}浏览器};
  \draw[flow] (host) -- (rt);
  \draw[flow] (rt)   -- (phone);
  \node[font=\scriptsize\sffamily\color{sddim}, below=0.15cm of host]  {http://192.168.x.x:8787};
\end{tikzpicture}
\end{center}

\textbf{特点。} 零延迟($\sim$\,30-80\,ms 端到端),不依赖外网,完全自治。一旦路由器和导播机连上同一网段,主持人扫的二维码就是导播工作机自身的局域网 IP + token。

\begin{warnbox}
\textbf{现场坑两条}:

(1)\,\textbf{客户端隔离}。某些场地(尤其大型展馆 / 教学楼)的 WiFi 上做了 AP isolation,同一 WiFi 里两台设备互相看不见,LAN 模式直接走不通。测试当天务必带备用便携路由器或手机热点,在策划会上和场地方明确确认这一条。

(2)\,\textbf{物理覆盖不到}。即使场地 WiFi 没做隔离,演播室到舞台十几米加墙加各种 2.4 GHz 干扰源(调音台、灯光控制器、其他无线话筒、观众的手机),一台便携路由器并不一定撑得住这段距离。Ping 通和"翻页延迟可接受"之间还有几个数量级的差距。预演要在演出地点的真实位置测试,不是会场入口或后台办公室。信号过不去就直接走 Relay,这是 Relay 模式存在的第二个理由。
\end{warnbox}

\subsection{Relay 模式 —— 跨网兜底}

\textbf{典型场景}:四种情况会让 LAN 走不通,Relay 接管:(1) 现场 WiFi 你不能配,场地方不给你接路由;(2) WiFi 做了客户端隔离;(3) 演播室到舞台距离 + 中间的墙和电子设备让一台便携路由根本覆盖不过去(这条是漫展场地最常见的真实问题);(4) 导播机和舞台在两栋楼。这四种情况里,只要导播机能上校园网、主持人手机能上校园网或外网,Relay 就能接管。

\begin{center}
\begin{tikzpicture}[
  node distance=1.2cm,
  box/.style={rectangle, draw=sdink, line width=0.7pt, rounded corners=2pt, minimum width=2.6cm, minimum height=0.95cm, align=center, font=\small\sffamily},
  relay/.style={rectangle, draw=sdaccent, line width=1pt, rounded corners=2pt, minimum width=2.6cm, minimum height=0.95cm, align=center, font=\small\sffamily, fill=sdpanel},
  flow/.style={-{Stealth[length=4pt]}, line width=0.7pt}
]
  \node[box]                      (host)  {导播工作机};
  \node[relay, right=of host]     (relay) {Relay 服务器\\\scriptsize\color{sddim}CLAB · 公网 IP};
  \node[box, right=of relay]      (phone) {主持人手机};
  \draw[flow] (host)  -- node[above, font=\scriptsize\sffamily\color{sddim}, align=center]{WSS 控制\\+ 二进制推流} (relay);
  \draw[flow] (relay) -- node[above, font=\scriptsize\sffamily\color{sddim}, align=center]{HTTPS\\+ MJPEG} (phone);
\end{tikzpicture}
\end{center}

\textbf{怎么工作。} 导播工作机主动开两条出站 WebSocket 到 Relay:控制隧道 + 二进制推流隧道。Relay 在公网开 80/8080,把主持人的 HTTP 请求经控制隧道反向打到导播机,把推流转写成普通 MJPEG 流给主持人浏览器。\textbf{业务路由代码两边完全一样} —— 导播侧的 \texttt{axum::Router} 被 Relay 通过 \texttt{tower::ServiceExt::oneshot} 当成 Service 调一次,响应再原路返回。主持人 URL 只是从 LAN 的 \texttt{:8787/presenter} 变成 Relay 的 \texttt{/r/<sid>/presenter},一行业务代码不用改。

\textbf{什么时候用、什么时候别用}:校园内跨楼、彩排时主持人不在现场提前演练、合作社团之间远程协调,这些场合 Relay 都合适。但对外公开演出暂时别走 Relay —— 写白皮书的时候,TLS、限速、抗 DDoS 这套都还在路线图上,Relay 仅在校园网内验证过。

设计上详见 ADR-0007(\texttt{docs/adr/0007-relay-hybrid-tunnel.md})。

\subsection{两种模式如何选}

\begin{center}\small
\begin{tabularx}{\linewidth}{l X X}
\toprule
\textbf{条件} & \textbf{LAN 模式} & \textbf{Relay 模式}\\
\midrule
现场能控 WiFi(自带路由器 / 热点) & \cellcolor{sdpanel}\bfseries 首选 & 备选 \\
场地 WiFi 客户端隔离 & 失效 & \cellcolor{sdpanel}\bfseries 唯一选择 \\
导播机和主持人分两栋楼 & 通常失效 & \cellcolor{sdpanel}\bfseries 首选 \\
延迟敏感(走位拍点) & \cellcolor{sdpanel}\bfseries 优 ($\sim$50\,ms) & 一般 ($\sim$80-200\,ms) \\
依赖外网可达性 & 否 & \cellcolor{sdred!10}\bfseries 是 \\
\bottomrule
\end{tabularx}
\end{center}

怎么搭配?LAN 为主,Relay 当兜底。导播 GUI 里把 Relay 配置常驻着,关键演出当天保持 Relay 状态绿灯;真到 LAN 抽风的时刻,切二维码就完事,30 秒内。
